\BookPart{Programming molecular dynamics}{dynamics}{%
Since about 2006, DNA computing has been transformed by an enzyme-free
mechanism: toehold-mediated strand displacement. Because one strand can
displace another at a rate set by the length of a short sequence handle, a
designer can program reactions directly in sequence. Part~V shows how this
single mechanism yields chemical reaction networks that behave like programs,
digital logic circuits with hundreds of components, analog computation and DNA
neural networks that classify and learn, and, through self-assembly, crystals
that execute algorithms and origami that folds into nanoscale shapes.

\textbf{Chapters in this part}
\begin{itemize}\setlength{\itemsep}{2pt}
\item \textbf{19\quad Toehold-mediated strand displacement.} How does an invading strand replace an incumbent, and what controls the rate?
\item \textbf{20\quad Chemical reaction networks as programs.} How can a set of reactions compute a function, and how is it compiled into DNA?
\item \textbf{21\quad Digital, analog and neural DNA circuits.} How are thresholds, restoration and weighted sums built, and how have DNA neural networks learned?
\item \textbf{22\quad Self-assembly, tiles and geometry.} How can growth itself compute, and how does a scaffold fold into a designed shape?
\end{itemize}}
